\documentclass[11pt,a4paper]{ctexart}

% 小标题不居中，左对齐
\ctexset{
    section/format = {\centering\bfseries\large},
    subsection/format = {\raggedright\bfseries\normalsize},
}

\usepackage{amsmath,amssymb}
\usepackage{graphicx}
\usepackage{booktabs}
\usepackage[margin=1in]{geometry}
\usepackage[hidelinks]{hyperref}
\usepackage{xcolor}

\title{EDAP 原型验证实验总结}
\author{}
\date{2026-08-12}

\begin{document}
\maketitle

\begin{center}
\small\href{https://github.com/Helios1919/EDAP-External-Depth-Attention-Plugin/tree/exp/edap-poc}{github.com/Helios1919/EDAP-External-Depth-Attention-Plugin}
\end{center}

\section{问题背景}

大模型在 RAG（Retrieval-Augmented Generation，检索增强生成）场景下经常出现知识冲突（knowledge conflict）：外部检索到的上下文和模型自己记住的参数化知识（parametric knowledge）不一致时，模型倾向于相信自己的记忆，导致幻觉（hallucination）。已有的方法如 CAD（Context-Aware Decoding，上下文感知解码，logit 层面做对比）和 DoLa（Decoding by Contrasting Layers，层间对比解码，不同层之间做对比）都是固定规则，不能针对特定数据分布做适配。

我想验证一个简单的想法：在冻结的大模型不同深度插入可训练的注意力插件，让插件自己学会「该信任浅层（上下文，context）还是深层（参数知识，parametric knowledge）」——这就是 EDAP（External Depth Attention Plugin，外部深度注意力插件）。

\section{实验设计与方法}

\subsection{模型架构}

\begin{itemize}
    \item 主干网络（backbone）：Qwen2.5-7B，完全冻结
    \item 28 层 Transformer 按 7 层一组切成 4 个块（block），块边界处插入 EDAP 插件
    \item 每个插件本质是一个小型多头交叉注意力（multi-head cross-attention，8 个头），Q 来自当前块的输出，K/V 来自所有之前块的输出（含嵌入层输出）
    \item 三层 LayerNorm（Layer Normalization，层归一化：输入/键/输出）+ 可学习的深度位置嵌入（depth position embedding）+ 零初始化 $\mathbf{W}_O$（保证初始时等价于恒等映射）
    \item LM Head（Language Modeling Head，语言建模头）解冻，否则梯度无法回传到插件
    \item 可训练参数：EDAP 206M + LM Head 545M = 751M（占主干 9.9\%）
\end{itemize}

\subsection{数据集}

\begin{itemize}
    \item \textbf{训练集}：ConFiQA，三种场景类型——counterfactual（反事实：上下文与参数化知识矛盾）、context\_required（需要上下文：参数化知识不足，必须使用上下文）、context\_irrelevant（上下文无关：参数化知识充足，上下文为干扰项）。反事实样本额外翻转为正确上下文变体（flipped variant），以教会模型区分场景而非学习固定偏置。数据增强后约 7000 条。
    \item \textbf{评估集}：ConFiQA（同分布，in-distribution）+ NQ-Swap（异分布，out-of-distribution）。NQ-Swap 是将 Natural Questions 段落中的命名实体进行替换制造的实体替换型反事实数据集，用于检测插件学到的是通用的「信任上下文」启发式规则，还是 ConFiQA 特定的路由策略（routing strategy）。
\end{itemize}

\subsection{训练配置}

\begin{table}[h]
\centering
\begin{tabular}{ll}
\toprule
参数 & 值 \\
\midrule
Epochs & 5（早停 patience=2） \\
Batch size & 8 $\times$ grad\_accum 2 = 有效 16 \\
学习率 & 1e-4，线性预热 450 步 $\to$ 余弦衰减 \\
优化器 & AdamW，weight\_decay=0.01 \\
标签平滑 & 0.1 \\
验证集比例 & 20\% \\
精度 & bfloat16（A100 40GB） \\
\bottomrule
\end{tabular}
\end{table}

损失函数基于全序列 teacher forcing（教师强制），在所有答案 token 上计算交叉熵（cross-entropy），仅对答案区间进行掩码。除 LM Head 外，主干网络所有参数均冻结。

\subsection{基线方法}

\begin{itemize}
    \item \textbf{Greedy（贪心解码）}：不做任何干预的标准贪心解码（greedy decoding），即冻结的 Qwen2.5-7B 主干原样输出。
    \item \textbf{CAD（Context-Aware Decoding）}：输出 logit 对比。每个解码步，logit 计算为 $(1+\alpha)\,\text{logits}_{\text{ctx}} - \alpha\,\text{logits}_{\text{no-ctx}}$，其中 $\alpha = 1.0$。
    \item \textbf{DoLa（Decoding by Contrasting Layers）}：内部层对比。将最终层（第 28 层）的 logit 与早期退出层（early exit，第 13 层）的 logit 进行对比，两者均通过最终 LayerNorm 和 LM Head。
\end{itemize}

所有方法均使用多 token 贪心生成（最多 32 个新 token），以确保公平比较。额外引入 P-EM（Prefix Exact Match，前缀精确匹配）指标，因为在实验中发现 CAD/DoLa 会抑制 EOS（End of Sequence，序列结束标记），导致模型生成正确答案后无法停止。

\section{实验结果}

\subsection{总体结果}

\begin{table}[h]
\centering
\begin{tabular}{lccc}
\toprule
方法 & ConFiQA EM & NQ-Swap EM & NQ-Swap P-EM \\
\midrule
Greedy & 16.60\% & 44.00\% & 63.60\% \\
CAD   & 3.60\%  & 7.40\%  & 64.80\% \\
DoLa  & 5.40\%  & 7.80\%  & 23.20\% \\
\textbf{EDAP} & \textbf{58.20\%} & 10.80\% & 14.80\% \\
\bottomrule
\end{tabular}
\end{table}

\subsection{ConFiQA 按答案来源（correct\_source）分解}

\begin{table}[h]
\centering
\begin{tabular}{lcc}
\toprule
方法 & 上下文型 (n=245) & 记忆型 (n=255) \\
\midrule
Greedy & 28.7\% & 3.7\% \\
CAD   & 7.0\%  & 0.0\% \\
DoLa  & 9.7\%  & 0.8\% \\
\textbf{EDAP} & \textbf{77.1\%} & \textbf{38.0\%} \\
\bottomrule
\end{tabular}
\end{table}

\section{分析与讨论}

\begin{enumerate}
    \item \textbf{EDAP 在 ConFiQA 上效果显著。} ConFiQA 从 16.6\% EM 提升至 58.2\%（+41.6pp），上下文型（correct\_source = context）达 77.1\%。说明插件确实学会了区分不同冲突类型，且对「需要依赖上下文」的场景几乎能做到可靠跟随。

    \item \textbf{不迁移到 NQ-Swap 是正向信号。} EDAP 在 NQ-Swap 上仅获 10.8\% EM，比 Greedy 的 44.0\% 还低。说明插件\textit{没有}学到简单的「多信任上下文」启发式——如果学到了，NQ-Swap（上下文始终为正确答案）也应该受益。这个退化（degradation）证明路由策略是 ConFiQA 特定的，支持模块化思路：不同分布的冲突可以各自训练专用插件，组合使用而互不干扰。

    \item \textbf{CAD 的 EOS 抑制问题得到定量验证。} CAD 在 NQ-Swap 上的 P-EM 高达 64.8\%，但 EM 仅 7.4\%，差距 57.4pp——模型实际上知道正确答案，但 logit 对比操作 $(1+\alpha)\text{logits}_{\text{ctx}} - \alpha\text{logits}_{\text{no-ctx}}$ 抑制了 EOS token，导致解码无法正常终止。这验证了多 token 生成范式以及 P-EM 作为诊断指标的必要性。

    \item \textbf{Greedy 解码存在显著的领域敏感性。} 冻结 Qwen2.5-7B 在 NQ-Swap 上可达 44.0\%（实体替换，本质是阅读理解），但在 ConFiQA 上仅 16.6\%（需要模型判断「我是否知道这件事」）。这种不对称性表明原生 LLM 擅长从上下文中提取事实，但在面对需要元认知判断的微妙知识冲突时能力很弱。
\end{enumerate}

\section{总结与后续计划}

这个原型实验验证了 EDAP 思路可行：在冻结 LLM 上插入可训练的跨深度注意力模块，让模型学到针对特定分布的冲突消解策略（conflict resolution strategy）。ConFiQA 上效果显著（+41.6pp EM），NQ-Swap 上的退化反而是好事——说明学到的不是粗糙启发式，而是分布特定的路由。

后续计划：
\begin{itemize}
    \item 跑 shuffle-depth 消融实验（ablation study），随机打乱深度顺序，直接量化深度维度对性能的贡献
    \item 尝试为不同冲突类型分别训练独立的 EDAP 实例，然后组合使用（composition）
    \item 换其他主干架构（architecture）和模型规模（scale），验证方法的通用性
\end{itemize}

\end{document}
